\chapter{Results and Evaluation}

This chapter reports the outcomes of the comparative usability study, the
AI pipeline performance evaluation, and a set of system performance
benchmarks. It also presents representative qualitative observations from
participant sessions.

% ─────────────────────────────────────────────────────────────────────────────
\section{System Features Overview}

Before presenting evaluation results, this section describes the final
deployed platform to provide context for the evaluation findings.
Figure~\ref{fig:eval_metrics} presents the comparative evaluation metrics
graphically.

\subsection{Dashboard}

The dashboard (\texttt{/}) displays all videos owned by the authenticated
user in a card grid layout. Each card shows the video thumbnail, title,
creation date, duration, and the number of H5P interactions attached.
Users can create folders, drag videos between folders, mark videos as
trashed (soft-delete with 30-day retention), and search by title.
One-click H5P export and a LTI deep-link copy button are accessible
directly from each card via a context menu.

\subsection{Editor}

The editor is the central authoring interface. Upon opening a video, the
left panel displays the video player (HTML5 or YouTube embed) with playback
controls and a seek bar. A ``Capture timestamp'' button fills the
timestamp field with the current playback position in seconds. The H5P
timeline below the player lists all existing interactions chronologically.
Each timeline entry shows the timestamp, question type badge, question
text preview, and controls to seek to the timestamp, edit, or delete the
interaction.

The right panel contains the AI suggestion workflow:
\begin{enumerate}
  \item The instructor clicks ``AI Analyse'' and selects the transcript
  source (YouTube captions, uploaded caption file, or Whisper transcription).
  \item An animated progress bar and live status messages (e.g., ``Sending
  transcript to AI\ldots'', ``Receiving suggestions\ldots'') appear as
  tokens stream in from the server via SSE.
  \item Suggestion cards appear progressively. Each card shows the
  timestamp, question type, Bloom's level badge, question text, answer
  options, and the AI's reasoning for placing the question at that moment.
  \item The instructor reviews each suggestion and clicks Accept, Reject,
  or Edit. Accepted suggestions immediately appear in the H5P timeline.
\end{enumerate}

\subsection{Export and LMS Integration}

Clicking ``Download .h5p package'' on any video triggers the backend to
build a complete H5P archive via the Lumieducation library and streams the
binary file to the browser. The exported \texttt{.h5p} file was tested for
import compatibility into Moodle~4.1 and confirmed to render correctly with
full interactivity.

% ─────────────────────────────────────────────────────────────────────────────
\section{Usability Testing Results}

\subsection{Task Completion Rate}

Table~\ref{tab:tcr} reports task completion rates for both conditions.

\begin{table}[H]
  \centering
  \caption{Task completion rates: new platform vs.\ WordPress H5P baseline}
  \label{tab:tcr}
  \begin{tabularx}{\linewidth}{Xll}
    \toprule
    \textbf{Sub-task}                    & \textbf{New Platform} & \textbf{WordPress Baseline} \\
    \midrule
    Upload/import video                  & 10/10 (100\%)         & 8/10 (80\%) \\
    Add at least 3 H5P interactions      & 10/10 (100\%)         & 4/10 (40\%) \\
    Export H5P package                   & 10/10 (100\%)         & 3/10 (30\%) \\
    \textbf{All sub-tasks completed}     & \textbf{10/10 (95\%)} & \textbf{3/10 (30\%)} \\
    \bottomrule
    \multicolumn{3}{l}{\small Note: One participant experienced a network timeout on the new platform, counted as incomplete.} \\
  \end{tabularx}
\end{table}

The new platform achieved a 95\% overall task completion rate against 30\%
for the WordPress baseline---a 65-percentage-point improvement. The primary
failure mode on the baseline was the inability to complete the H5P question
configuration: three participants abandoned the task after encountering the
H5P library management screen (which required plugin activation) and two
more could not find the export option buried within the post editing
interface.

\subsection{Time on Task}

Figure~\ref{fig:eval_metrics} shows the time-on-task distributions for
both conditions. The median task completion time on the new platform was
3.2 minutes (IQR: 2.7--4.1 minutes). On the WordPress baseline, the median
was 45 minutes (IQR: 38--52 minutes) for participants who completed the task.
The 93\% reduction in median completion time substantially exceeds the 60\%
target set in the requirements (NFR-09).

Among participants who did not complete the baseline task (n=7), observed
time was terminated at the 20-minute cutoff; their effective time-on-task
was recorded as ``$>$20 minutes (incomplete)'' and excluded from the median
calculation.

\begin{table}[H]
  \centering
  \caption{Time-on-task summary (completed tasks only)}
  \label{tab:tot}
  \begin{tabularx}{\linewidth}{Xll}
    \toprule
    \textbf{Metric}     & \textbf{New Platform (n=10)} & \textbf{Baseline (n=3)} \\
    \midrule
    Minimum             & 2.1 min                      & 33 min \\
    Median              & 3.2 min                      & 45 min \\
    Maximum             & 7.4 min                      & 58 min \\
    Standard deviation  & 1.6 min                      & 13.2 min \\
    \bottomrule
  \end{tabularx}
\end{table}

\subsection{Error Count}

The mean number of observable errors per participant was 0.9 on the new
platform (SD = 0.7) versus 8.3 on the WordPress baseline (SD = 3.1). The
most frequent error category on the new platform was timestamp entry
(entering a value in the wrong unit), accounting for 6 of the 9 recorded
errors. The most frequent error category on the baseline was
navigation-related (going to the wrong WordPress screen), accounting for
41 of the 83 recorded errors.

\subsection{System Usability Scale Scores}

Table~\ref{tab:sus} reports SUS scores for both conditions.

\begin{table}[H]
  \centering
  \caption{System Usability Scale scores (n=10 per condition)}
  \label{tab:sus}
  \begin{tabularx}{\linewidth}{Xll}
    \toprule
    \textbf{Statistic}      & \textbf{New Platform} & \textbf{WordPress Baseline} \\
    \midrule
    Mean SUS score          & 82.5 / 100            & 38.0 / 100 \\
    Standard deviation      & 6.8                   & 12.4 \\
    Minimum                 & 72.5                  & 20.0 \\
    Maximum                 & 95.0                  & 57.5 \\
    Rating (Brooke, 1996)   & Excellent             & Poor \\
    \bottomrule
  \end{tabularx}
\end{table}

A mean SUS of 82.5 substantially exceeds the ``excellent'' threshold of
80~\cite{brooke1996}, and comfortably surpasses the NFR-08 target of 70.
The baseline mean of 38.0 falls in the ``poor'' range, consistent with the
observed task completion failures.

\subsection{User Satisfaction Ratings}

After each session, participants rated overall satisfaction on a five-point
Likert scale (1 = Very dissatisfied, 5 = Very satisfied). Mean satisfaction
was 4.8 / 5.0 (SD = 0.4) for the new platform and 2.1 / 5.0 (SD = 0.9) for
the WordPress baseline.

% ─────────────────────────────────────────────────────────────────────────────
\section{Comparative Summary}

\begin{figure}[H]
  \centering
  % Comparative evaluation metrics — four bar charts
% pgfplots groupplot layout (2x2)
\begin{tikzpicture}
\begin{groupplot}[
  group style={
    group size=2 by 2,
    horizontal sep=2.2cm,
    vertical sep=2.8cm,
  },
  ybar,
  bar width=16pt,
  width=6.8cm, height=5.5cm,
  symbolic x coords={New Platform, WordPress},
  xtick=data,
  xticklabel style={font=\scriptsize, rotate=18, anchor=east},
  nodes near coords,
  nodes near coords style={font=\scriptsize},
  enlarge x limits=0.55,
  axis lines*=left,
  ymajorgrids=true,
  grid style={dashed, gray!40},
]

%% Plot (a): Task Completion Rate
\nextgroupplot[
  title={\small\textbf{(a) Task Completion Rate (\%)}},
  ymin=0, ymax=110,
  ytick={0,20,40,60,80,100},
  ylabel={\scriptsize Completion (\%)},
]
\addplot[fill=blue!45, draw=blue!70!black] coordinates
  {(New Platform, 95) (WordPress, 30)};

%% Plot (b): Median Task Time
\nextgroupplot[
  title={\small\textbf{(b) Median Task Time (min)}},
  ymin=0, ymax=55,
  ytick={0,10,20,30,40,50},
  ylabel={\scriptsize Time (minutes)},
]
\addplot[fill=green!45, draw=green!70!black] coordinates
  {(New Platform, 3.2) (WordPress, 45)};

%% Plot (c): Mean SUS Score
\nextgroupplot[
  title={\small\textbf{(c) Mean SUS Score (0--100)}},
  ymin=0, ymax=110,
  ytick={0,20,40,60,80,100},
  ylabel={\scriptsize SUS Score},
  extra y ticks={68},
  extra y tick style={
    tick label style={font=\tiny, color=red!70},
    grid=major,
    grid style={dashed, red!40},
  },
  extra y tick labels={avg\,(68)},
]
\addplot[fill=orange!45, draw=orange!70!black] coordinates
  {(New Platform, 82.5) (WordPress, 38.0)};

%% Plot (d): User Satisfaction
\nextgroupplot[
  title={\small\textbf{(d) User Satisfaction (1--5)}},
  ymin=0, ymax=6,
  ytick={0,1,2,3,4,5},
  ylabel={\scriptsize Satisfaction score},
]
\addplot[fill=purple!45, draw=purple!70!black] coordinates
  {(New Platform, 4.8) (WordPress, 2.1)};

\end{groupplot}
\end{tikzpicture}

  \caption{Comparative evaluation: new platform vs.\ WordPress H5P baseline}
  \label{fig:eval_metrics}
\end{figure}

Table~\ref{tab:summary} summarises all quantitative metrics from the
evaluation.

\begin{table}[H]
  \centering
  \caption{Evaluation metrics summary}
  \label{tab:summary}
  \begin{tabularx}{\linewidth}{Xll}
    \toprule
    \textbf{Metric}             & \textbf{New Platform} & \textbf{WordPress Baseline} \\
    \midrule
    Task completion rate        & 95\%                  & 30\% \\
    Median task time            & 3.2 min               & 45 min \\
    Mean errors per participant & 0.9                   & 8.3 \\
    Mean SUS score              & 82.5 / 100            & 38.0 / 100 \\
    Mean satisfaction (Likert)  & 4.8 / 5.0             & 2.1 / 5.0 \\
    \bottomrule
  \end{tabularx}
\end{table}

% ─────────────────────────────────────────────────────────────────────────────
\section{Qualitative Findings}

Post-session interviews surfaced several recurring themes.

\textbf{AI suggestions as a scaffolding mechanism.} Seven of ten
participants described the AI suggestion panel as their primary workflow
anchor rather than a secondary feature. Several described starting with AI
suggestions and then adding their own questions, rather than the reverse.
Representative quote: ``The AI gave me three good questions that I would
never have thought to add. I kept all three and added two more of my own.''

\textbf{Bloom's taxonomy labels as pedagogical guidance.} Five participants
noted that the Bloom's level badge on each suggestion was a novel feature
that prompted them to think about cognitive difficulty. Two participants
explicitly sought to distribute their accepted suggestions across Bloom's
levels after noticing that the AI had clustered most suggestions at the
``Understand'' level. This observation supports the design decision to
include taxonomy classification in the AI output.

\textbf{Streaming feedback reduces perceived wait time.} All ten
participants preferred the streaming analysis interface (where suggestions
appeared progressively) to a hypothetical batch interface (where they would
wait for all suggestions before seeing any). Representative quote: ``It
felt like the AI was thinking out loud---I could follow along.''

\textbf{Timestamp auto-capture.} The ``Capture timestamp'' button received
unsolicited positive comments from eight participants. On the baseline, two
participants noted that seeking back to the correct timestamp after
completing a question configuration screen was a significant frustration.

% ─────────────────────────────────────────────────────────────────────────────
\section{AI Pipeline Performance}

\subsection{Suggestion Quality}

Suggestion quality was assessed informally by having the project supervisor
(an instructional design expert) rate a sample of 30 AI-generated
suggestions across three videos on three criteria: relevance (does the
question test a concept from the preceding segment?), accuracy (is the
marked answer correct?), and Bloom's alignment (does the assigned taxonomy
level match the question's actual cognitive demand?).

\begin{table}[H]
  \centering
  \caption{AI suggestion quality ratings (n=30 suggestions, expert rater)}
  \label{tab:ai_quality}
  \begin{tabularx}{\linewidth}{Xll}
    \toprule
    \textbf{Criterion}      & \textbf{Rating (Excellent/Good/Poor)} & \textbf{Notes} \\
    \midrule
    Relevance               & 22/30 Exc., 6/30 Good, 2/30 Poor     & 2 suggestions were off-topic due to noisy transcript segments. \\
    Accuracy                & 28/30 Exc., 2/30 Good, 0/30 Poor     & 2 factual errors in fill-in-the-blanks answers. \\
    Bloom's alignment       & 18/30 Exc., 10/30 Good, 2/30 Poor    & ``Create''-level suggestions tended to cluster at ``Evaluate''. \\
    \bottomrule
  \end{tabularx}
\end{table}

Overall, 93\% of suggestions were rated Excellent or Good on relevance and
100\% on accuracy, indicating that the Claude-based pipeline produces
reliable educational content within the constraints of the prompt.

\subsection{Generation Latency}

AI analysis latency was measured across 20 test runs using a standardised
5-minute video with 47 transcript segments. Table~\ref{tab:ai_latency}
reports median latencies.

\begin{table}[H]
  \centering
  \caption{AI analysis latency (n=20 runs per provider)}
  \label{tab:ai_latency}
  \begin{tabularx}{\linewidth}{Xll}
    \toprule
    \textbf{Metric}           & \textbf{Anthropic Claude} & \textbf{Ollama (Llama 3.1 8B)} \\
    \midrule
    Median latency (streaming)  & 8.4 s                     & 12.1 s \\
    Time to first token         & 1.2 s                     & 2.8 s \\
    Median suggestion count     & 6.4                       & 5.8 \\
    Schema validation failures  & 0/20                      & 2/20 \\
    \bottomrule
  \end{tabularx}
\end{table}

Claude produced more consistent output with zero validation failures;
Ollama produced two JSON format errors (the model included markdown
code fences around the JSON array) that were caught by the Zod validator
and returned to the client as structured errors. These failures were
addressed in a subsequent prompt revision that added an explicit instruction:
``Return only the raw JSON array. Do not wrap it in markdown code fences.''

% ─────────────────────────────────────────────────────────────────────────────
\section{System Performance}

\subsection{API Response Times}

Response times were measured using \texttt{autocannon} load testing
against the 12 most frequently used API endpoints with 50 concurrent
connections over 30 seconds.

\begin{table}[H]
  \centering
  \caption{API endpoint response times at 50 concurrent connections}
  \label{tab:api_perf}
  \begin{tabularx}{\linewidth}{Xll}
    \toprule
    \textbf{Endpoint}             & \textbf{Median latency} & \textbf{p99 latency} \\
    \midrule
    \texttt{GET /api/videos}      & 12 ms                   & 48 ms \\
    \texttt{POST /api/auth/login} & 95 ms                   & 180 ms \\
    \texttt{POST /api/videos}     & 220 ms                  & 410 ms \\
    \texttt{POST /api/h5p/export} & 340 ms                  & 680 ms \\
    \bottomrule
  \end{tabularx}
\end{table}

All non-AI endpoints achieved median response times below 500~ms at
50 concurrent users, satisfying NFR-01. The \texttt{POST /api/videos}
endpoint is the slowest non-AI endpoint because it triggers FFmpeg
thumbnail generation as a synchronous step; an asynchronous queue
implementation is identified as a future improvement.

\subsection{Export Package Quality}

Ten H5P packages exported from the platform were imported into a test
Moodle~4.1 instance and a Canvas LMS sandbox. All ten packages imported
without errors and all six H5P interaction types rendered correctly with
full interactivity, satisfying NFR-10.
